\cotefile{157-5}
\scanpage{373}
\mspage{289}
% archive sheet number: 372

\uline{4.13.5.} Soit $\underline S$ une sous-catégorie \uuline{str. pleine}
de $M$ (catégorie accessible), \inserted{$\pi'_0$ un cardinal infini.}
Dire que $\underline S$ est $\pi'_0$-accessible dans $M$ signifie que
les conditions suivantes sont satisfaites

1) $\pi'_0$ est utile pour $\underline S$, i.e. $\underline S \overset{\sim}{\underset{\sim}{\leftarrow}}\operatorname{Ind}_{\pi'_0}(S_{\pi'_0})$\unsure{;}

2) $\underline S\hookrightarrow M$ est $\pi'_0$-accessible

3) $\underline S_{\pi'_0}\subset M_{\pi'_0}$ (qui implique que $S_{\pi'_0}$ est
équivalente à une petite catégorie).

\deleted{Soit} Supposons ces conditions satisfaites,
et soit $\pi^{\unsure{\prime|\diamond}}$ un cardinal $>\pi'_0$. Est-ce que
$\underline S$ est $\pi^{\unsure{\prime}}$-accessible \inserted{dans $M$}, i.e. les conditions
similaires sont-elles satisfaites pour $\pi$?
C'est évident pour $2^\circ$), mais risque d'être
faux pour $1^\circ$, $3^\circ$.

\deleted{\deleted{C'est pourquoi il a été utile \unsure{d'introduire|de déterminer}} le plus petit cardinal
$\pi'_0$ \inserted{\unsure{infini}} tel que $\pi'_0$ soit utile pour $\underline S$, et
que $\underline S\hookrightarrow M$ soit $\pi'_0$-accessible, et
$\pi'_1$ \deleted{$=\operatorname{Sup}(\pi'_0,\operatorname{card}\unsure{\operatorname{tot}|\operatorname{Fl}}S_{\pi'_0})$} le plus petit
cardinal $\geq$\inserted{$\pi'_0$,} $(\pi'_0,\operatorname{card}\operatorname{tot}S_{\pi'_0}$, et tel que
$S_{\pi'_0}\subset\deleted{\unsure{\cong}}M_{\pi'_1}$. \deleted{Si alors} (NB $\pi_0$ est le plus
petit cardinal \deleted{tel que} utile pour $M$.)}
\deleted{Si alors $\pi\geq\pi'_1$} Mais faisons appel à
% slanted
\margin{\deleted{\unsure{En fait, distinguons} \unsure{rel.} \unsure{à} \illegible{1} \unsure{4.13.1} à \unsure{4.13.2} \unsure{ici}.} \unsure{4} \unsure{4.}}
4.13.2. \unsure{c)}. \deleted{Supposons que $\pi'_0\geq\pi_0$,} Posons
\[
\pi'_1=\operatorname{Sup}(\pi'_0,\operatorname{card}\operatorname{tot} S_{\pi'_0},\pi_0)
\]
\margin{\unsure{$\Vert$|``}plus petit cardinal infini \deleted{\unsure{tel que}} utile pour $M$\unsure{.}}
